\documentclass[10pt,a4paper]{amsart}
\usepackage{amsmath,amssymb,amsthm,graphicx}
% Retained minimal preamble for the Stacks Project chapter more-algebra.tex (More on Algebra).
% The chapter file pinned-original-source.tex carries only \input{preamble} on its first line, so no
% preamble text lives inside the chapter itself.  The Stacks Project uses one single shared preamble for
% every chapter; the snapshot stages it identically as sources/stacks-chow/preamble.tex and
% sources/stacks-etale-cohomology/preamble.tex.  The block below is copied mechanically, line for line,
% from sources/stacks-chow/preamble.tex lines 14-18 (the Xy-pic setup) and 154-217 (the native theorem
% declarations together with the mathematical macros).  Nothing was added, removed or redefined; the
% wrapper class and the amsmath/amssymb/amsthm/graphicx base packages are supplied by the assembler.
% Chapter: more-algebra.tex -- More on Algebra.
% For commutative diagrams we use Xy-pic
\usepackage[all]{xy}

% We use 2cell for 2-commutative diagrams.
\xyoption{2cell}

% Theorem environments.
%
\theoremstyle{plain}
\newtheorem{theorem}[subsection]{Theorem}
\newtheorem{proposition}[subsection]{Proposition}
\newtheorem{lemma}[subsection]{Lemma}

\theoremstyle{definition}
\newtheorem{definition}[subsection]{Definition}
\newtheorem{example}[subsection]{Example}
\newtheorem{exercise}[subsection]{Exercise}
\newtheorem{situation}[subsection]{Situation}

\theoremstyle{remark}
\newtheorem{remark}[subsection]{Remark}
\newtheorem{remarks}[subsection]{Remarks}

\numberwithin{equation}{subsection}

% Macros
%
\def\lim{\mathop{\mathrm{lim}}\nolimits}
\def\colim{\mathop{\mathrm{colim}}\nolimits}
\def\Spec{\mathop{\mathrm{Spec}}}
\def\Hom{\mathop{\mathrm{Hom}}\nolimits}
\def\Ext{\mathop{\mathrm{Ext}}\nolimits}
\def\SheafHom{\mathop{\mathcal{H}\!\mathit{om}}\nolimits}
\def\SheafExt{\mathop{\mathcal{E}\!\mathit{xt}}\nolimits}
\def\Sch{\mathit{Sch}}
\def\Mor{\mathop{\mathrm{Mor}}\nolimits}
\def\Ob{\mathop{\mathrm{Ob}}\nolimits}
\def\Sh{\mathop{\mathit{Sh}}\nolimits}
\def\NL{\mathop{N\!L}\nolimits}
\def\CH{\mathop{\mathrm{CH}}\nolimits}
\def\proetale{{pro\text{-}\acute{e}tale}}
\def\etale{{\acute{e}tale}}
\def\QCoh{\mathit{QCoh}}
\def\Ker{\mathop{\mathrm{Ker}}}
\def\Im{\mathop{\mathrm{Im}}}
\def\Coker{\mathop{\mathrm{Coker}}}
\def\Coim{\mathop{\mathrm{Coim}}}

% Boxtimes
%
\DeclareMathSymbol{\boxtimes}{\mathbin}{AMSa}{"02}

%
% Macros for moduli stacks/spaces
%
\def\QCohstack{\mathcal{QC}\!\mathit{oh}}
\def\Cohstack{\mathcal{C}\!\mathit{oh}}
\def\Spacesstack{\mathcal{S}\!\mathit{paces}}
\def\Quotfunctor{\mathrm{Quot}}
\def\Hilbfunctor{\mathrm{Hilb}}
\def\Curvesstack{\mathcal{C}\!\mathit{urves}}
\def\Polarizedstack{\mathcal{P}\!\mathit{olarized}}
\def\Complexesstack{\mathcal{C}\!\mathit{omplexes}}
% \Pic is the operator that assigns to X its picard group, usage \Pic(X)
% \Picardstack_{X/B} denotes the Picard stack of X over B
% \Picardfunctor_{X/B} denotes the Picard functor of X over B
\def\Pic{\mathop{\mathrm{Pic}}\nolimits}
\def\Picardstack{\mathcal{P}\!\mathit{ic}}
\def\Picardfunctor{\mathrm{Pic}}
\def\Deformationcategory{\mathcal{D}\!\mathit{ef}}

% --- appended 2026-10-12 (Chief): the project own definitions of the reference/slogan/history
% display environments, copied mechanically from the shared Stacks preamble lines 8-12
% (\usepackage{verbatim} plus the three \newenvironment lines, which the chapter preamble keeps
% verbatim from the project preamble). Nothing added, removed or redefined: these environments
% carry no mathematical content and the project itself typesets them as comments. Without these
% three lines a counted statement whose own text uses them fails with Environment slogan undefined.
\usepackage{verbatim}
\newenvironment{reference}{\comment}{\endcomment}
\newenvironment{slogan}{\comment}{\endcomment}
\newenvironment{history}{\comment}{\endcomment}
\usepackage{cleveref}
\crefname{originallemma}{original lemma}{original lemmas}
\Crefname{originallemma}{Original lemma}{Original lemmas}
\makeatletter
\newlabel{lemma-henselization-noetherian}{{06LJ}{source}{}{sourcepointer-1}{}}
\newlabel{lemma-henselization-noetherian@cref}{{[originallemma][0][]06LJ}{[1][0][]source}{}{}{}}
\makeatother
\crefname{originallemma}{original lemma}{original lemmas}
\Crefname{originallemma}{Original lemma}{Original lemmas}
\makeatletter
\newlabel{algebra-lemma-flat-going-down}{{00HS}{source}{}{sourcepointer-2}{}}
\newlabel{algebra-lemma-flat-going-down@cref}{{[originallemma][0][]00HS}{[1][0][]source}{}{}{}}
\makeatother
\crefname{originallemma}{original lemma}{original lemmas}
\Crefname{originallemma}{Original lemma}{Original lemmas}
\makeatletter
\newlabel{algebra-lemma-injective-minimal-primes-in-image}{{00FK}{source}{}{sourcepointer-3}{}}
\newlabel{algebra-lemma-injective-minimal-primes-in-image@cref}{{[originallemma][0][]00FK}{[1][0][]source}{}{}{}}
\makeatother
\crefname{originaldefinition}{original definition}{original definitions}
\Crefname{originaldefinition}{Original definition}{Original definitions}
\makeatletter
\newlabel{definition-number-of-branches}{{0C26}{source}{}{sourcepointer-4}{}}
\newlabel{definition-number-of-branches@cref}{{[originaldefinition][0][]0C26}{[1][0][]source}{}{}{}}
\makeatother
\crefname{originallemma}{original lemma}{original lemmas}
\Crefname{originallemma}{Original lemma}{Original lemmas}
\makeatletter
\newlabel{algebra-lemma-complete-henselian}{{04GM}{source}{}{sourcepointer-5}{}}
\newlabel{algebra-lemma-complete-henselian@cref}{{[originallemma][0][]04GM}{[1][0][]source}{}{}{}}
\makeatother
\title{Branches of a Noetherian local ring versus its completion}
\author{The Stacks Project Authors}
\date{stacks-project:more-algebra.tex:local-snapshot-20261012; GNU FDL 1.2 or later}
\begin{document}
\maketitle
\setcounter{subsection}{1}
\noindent\textit{Editorial note.} All delivered mathematical text is the author's own, copied from the pinned Stacks Project chapter file more-algebra.tex; no statement or proof step was rewritten, reordered or removed. The counted span is the author environment and its own proof environment. References inside the span are handed to the assembler as navigation stubs that carry the canonical tag of the cited tag only; whether each of them is load-bearing is recorded on the stub itself.\par\medskip
\noindent\textit{Allowed standard prerequisites (editorial).}\begin{itemize}
\item The map from the henselization to the completion of a Noetherian local ring is flat and injective. invoked at proof lines 31949-31950. Sets up the minimal prime comparison.
\item Flat maps satisfy going down and injective ring maps realize every minimal prime of the target as lying over one from the source. invoked at proof lines 31950-31952. Proves part (1).
\item The completion is henselian and the number of branches is defined via minimal primes of henselizations. invoked at proof lines 31953-31956. Proves part (2).
\item The strict henselization of the completion is the base change of the strict henselization over the henselization. invoked at proof lines 31957-31958. Enables repeating the argument for part (3).
\end{itemize}\medskip
\ifdefined\hypertarget\hypertarget{sourcepointer-1}{}\fi
\noindent\textit{Source reference.} Original lemma 06LJ: sources/stacks-more-algebra/pinned-original-source.tex, source lines 10979--10979. Cited at proof lines 31949-31949 and 31957-31958 for the flat injective map and the strict henselization base change; the minimal prime and branch arguments are inside the counted span.\par\medskip
\ifdefined\hypertarget\hypertarget{sourcepointer-2}{}\fi
\noindent\textit{Source reference.} Original lemma 00HS: sources/stacks-chow/tags.txt, source lines 654--654. Cited at proof lines 31950-31951 as the first ingredient for part (1); the conclusion is inside the counted span.\par\medskip
\ifdefined\hypertarget\hypertarget{sourcepointer-3}{}\fi
\noindent\textit{Source reference.} Original lemma 00FK: sources/stacks-chow/tags.txt, source lines 574--574. Cited at proof lines 31951-31952 as the second ingredient for part (1); the conclusion is inside the counted span.\par\medskip
\ifdefined\hypertarget\hypertarget{sourcepointer-4}{}\fi
\noindent\textit{Source reference.} Original definition 0C26: sources/stacks-more-algebra/pinned-original-source.tex, source lines 31553--31553. Cited at proof lines 31954-31954 for the definitions behind parts (2) and (3); the inequalities are proved inside the counted span.\par\medskip
\ifdefined\hypertarget\hypertarget{sourcepointer-5}{}\fi
\noindent\textit{Source reference.} Original lemma 04GM: sources/stacks-chow/tags.txt, source lines 5838--5838. Cited at proof lines 31955-31956 to identify the branches of the completion; the inequality is inside the counted span.\par\medskip
\begin{lemma}
\label{lemma-nr-branches-completion}
Let $(A, \mathfrak m)$ be a Noetherian local ring.
\begin{enumerate}
\item The map $A^h \to A^\wedge$ defines a surjective map from minimal
primes of $A^\wedge$ to minimal primes of $A^h$.
\item The number of branches of $A$ is at most the number of branches
of $A^\wedge$.
\item The number of geometric branches of $A$ is at most the number
of geometric branches of $A^\wedge$.
\end{enumerate}
\end{lemma}

\begin{proof}
By Lemma \ref{lemma-henselization-noetherian} the map $A^h \to A^\wedge$
is flat and injective. Combining going down
(Algebra, Lemma \ref{algebra-lemma-flat-going-down}) and
Algebra, Lemma \ref{algebra-lemma-injective-minimal-primes-in-image}
we see that part (1) holds. Part (2) follows from this,
Definition \ref{definition-number-of-branches}, and
the fact that $A^\wedge$ is henselian
(Algebra, Lemma \ref{algebra-lemma-complete-henselian}).
By Lemma \ref{lemma-henselization-noetherian}
we have $(A^\wedge)^{sh} = A^{sh} \otimes_{A^h} A^\wedge$.
Thus we can repeat the arguments above using the flat injective
map $A^{sh} \to (A^\wedge)^{sh}$ to prove (3).
\end{proof}

\end{document}
